    \section{Correlation functions in QM (PSI QFT II, PS1)}
    
    Consider a 1-dimensional harmonic oscillator with Hamiltonian
    \begin{equation}
        \mathbf{H} = \frac{\mathbf{P}^2}{2m} + \frac{m\omega_0^2}{2}\mathbf{Q}^2
    \end{equation}
    We can write the propagator in real time
    \begin{equation}
        K(q_2, q_1; t) = \langle q_2 | U(t) | q_1 \rangle
    \end{equation}
    and the ``Euclidean'' propagator
    \begin{equation}
        K_E(q_2, q_1; \tau) = \langle q_2 | U(-i\tau) | q_1 \rangle
    \end{equation}
    using the time evolution operator
    \begin{equation}
        U(t) = \exp\!\left(\frac{t}{i\hbar}\mathbf{H}\right).
    \end{equation}
    We want to compute the ``Green's function'' or correlation function (also known as the propagator):
    \begin{equation}
        G(t_1, t_2) = \langle 0 | T(\mathbf{Q}(t_1)\mathbf{Q}(t_2)) | 0 \rangle
    \end{equation}
    $|0\rangle$ is the ground state of the harmonic oscillator and $T$ denotes the time-ordered product of the operators $\mathbf{Q}(t) = U(-t)\mathbf{Q}U(t)$ (in the Heisenberg picture). We start from the Euclidean Green's function at finite temperature
    \begin{equation}
        G_E(\tau_2, \tau_1; \beta) = \mathrm{Tr}\!\left(\boldsymbol{\rho}(\beta)\,\mathbf{Q}(-i\tau_2)\,\mathbf{Q}(-i\tau_1)\right)
    \end{equation}
    Here $\boldsymbol{\rho}(\beta)$ is the density matrix of the harmonic oscillator at temperature $T$ and $\beta = \frac{1}{k_B T}$.
    
    \begin{enumerate}[label=\alph*)]
    
        \item Why does this expression require $\tau_2 > \tau_1$?
        \begin{solution}
        We shall first demonstrate the equivalence between the real time and Euclidean time propagator. The real time propagator can be written as a path integral 
        \begin{equation}
            G(t_1,t_2) = \bra{0}T\mathbf{Q}(t_1)\mathbf{Q}(t_2)\ket{0} = \int\mathcal{D}q(t)q(t_1)q(t_2)e^{iS[q]/\hbar}.
        \end{equation}
        Since the time in the action $S[x]$ is merely an integration variable, we can perform a Wick rotation, and obtain
        \begin{equation}
            G(t_1,t_2) = \int\mathcal{D}q_E(\tau)q(-i\tau_1)q(-i\tau_2)e^{-S_E[q_E]/\hbar}.
        \end{equation}
       If we choose a periodic boundary condition, the propagator becomes an expectation value in the context of finite temperature statistical mechanics:
        \begin{equation}
            G_E(\tau_2, \tau_1;\beta) =\frac{1}{Z} \oint\mathcal{D}q_E(\tau)q(-i\tau_1)q(-i\tau_2)e^{-S_E[q_E]/\hbar} = \mathrm{Tr}\!\left(\boldsymbol{\rho}(\beta)\,\mathbf{Q}(-i\tau_2)\,\mathbf{Q}(-i\tau_1)\right).
        \end{equation}
        To see the above, we can write the path integral as 
        \begin{equation}
            \oint\mathcal{D}q_E(\tau) = \int dq_{E1}\int dq_{E2}\oint_{q_{Ei} = q_E(\tau_i)}\mathcal{D}q_E(\tau),
        \end{equation}
        using the completeness relation, we find
        \begin{align}
        &G_E(\tau_2, \tau_1;\beta ) \notag\\
        &= \frac{1}{Z}  \int dq_{E1}\int dq_{E2}
           \oint_{q_E(\tau_i)=q_{Ei}}\mathcal{D}q_E(\tau)\,
           q(-i\tau_1)q(-i\tau_2)\,e^{-S_E[q_E]/\hbar} \notag\\
           &= \frac{1}{Z} \int dq_{E1}\int dq_{E2}\int dq\, q_{E1}\,q_{E2}\,
           \langle q|e^{-(\beta-\tau_2)H/\hbar}|q_{E2}\rangle
           \langle q_{E2}|e^{-(\tau_2-\tau_1)H/\hbar}|q_{E1}\rangle
           \langle q_{E1}|e^{-\tau_1 H/\hbar}|q\rangle \notag\\
           &= \frac{1}{Z} \int dq_{E1}\int dq_{E2}\int dq\, 
           \langle q|e^{-(\beta-\tau_2)H/\hbar}\mathbf{Q}e^{-(\tau_2-\tau_1)H/\hbar}\mathbf{Q}e^{-\tau_1 H/\hbar}|q\rangle \notag\\
        &= \frac{1}{Z} \int dq\,
           \langle q|\,e^{-\beta H/\hbar}\,
           \mathbf{Q}(-i\tau_2)\,\mathbf{Q}(-i\tau_1)\,|q\rangle \notag\\
        &= \frac{1}{Z} \mathrm{Tr}\!\left(e^{-\beta H/\hbar}\,
           \mathbf{Q}(-i\tau_2)\,\mathbf{Q}(-i\tau_1)\right) \notag\\
        &= \mathrm{Tr}\!\left(\boldsymbol{\rho}(\beta)\,
           \mathbf{Q}(-i\tau_2)\,\mathbf{Q}(-i\tau_1)\right).
    \end{align}
            Expanding the trace over energy eigenstates $|n\rangle$ with $H|n\rangle = E_n|n\rangle$:
    \begin{align}
        G_E(\tau_2,\tau_1;\beta) 
        &= \mathrm{Tr}\!\left(\boldsymbol{\rho}(\beta)\,\mathbf{Q}(-i\tau_2)\,\mathbf{Q}(-i\tau_1)\right) \notag\\
        &= \frac{1}{Z}\sum_{m,n} e^{-\beta E_m}
           \langle m | e^{\tau_2 H}\,\mathbf{Q}\,e^{-\tau_2 H}
                       e^{\tau_1 H}\,\mathbf{Q}\,e^{-\tau_1 H} | m \rangle \notag\\
        &= \frac{1}{Z}\sum_{m,n} e^{-\beta E_m}\,e^{(\tau_2-\tau_1)(E_m - E_n)}
           |\langle m|\mathbf{Q}|n\rangle|^2
    \end{align}
    where $Z = \mathrm{Tr}(e^{-\beta H})$. For this sum to converge we need the exponential factor to be damped, which requires
    \begin{equation}
        \tau_2 - \tau_1 > 0 \qquad \Longrightarrow \qquad \tau_2 > \tau_1.
    \end{equation}
    If instead $\tau_2 < \tau_1$, the factor $e^{(\tau_2-\tau_1)(E_m-E_n)}$ diverges for large $E_m - E_n$, rendering the expression ill-defined. This is the Euclidean analogue of the time-ordering prescription in Minkowski space.
        \end{solution}
    
        \item Write $G_E(\tau_2, \tau_1; \beta)$ as a Euclidean functional integral.
        \begin{solution}
            As already demonstrated above, we have
            \begin{equation}
                G_E(\tau_2, \tau_1;\beta) =\frac{1}{Z} \oint\mathcal{D}q_E(\tau)q(-i\tau_1)q(-i\tau_2)e^{-S_E[q_E]/\hbar}
            \end{equation}
        \end{solution}
    
        \item Explain why $G_E = G_E(\tau, \beta)$ is only a function of $\tau = \tau_2 - \tau_1$, and why it should be considered as a periodic function in $\tau$. What is this period $\tau_\beta$ and what is its limit when the temperature $T$ is taken to be zero?
        \begin{solution}
            As demonstrated above, $G_E$, can be written as a function of $\tau = \tau_2-\tau_1$. This is due to the fact that the Hamiltonian is invariant under time translation. 
    
            The periodic boundary condition of the path integral can also be viewed as the Euclidean time itself being periodic. Therefore, we identify $\tau$ as $\tau + \tau_\beta$ with $\tau_\beta = \hbar \beta$. As $T\to 0$, $\tau_\beta \to \infty$.
        \end{solution}
    
        \item The Green's function $G_E(\tau; \beta)$ obeys the differential equation
        \begin{equation}
            \left(-\frac{\partial^2}{\partial\tau^2} + \omega_0^2\right) G_E(\tau;\beta) = \frac{\hbar}{m}\,\delta(\tau)
        \end{equation}
        Show that at zero temperature ($T = 0$, $\beta = \infty$), $G_E(\tau;\infty) = G_E(\tau)$ is given simply by
        \begin{equation}
            G_E(\tau) = \frac{\hbar}{m}\,\frac{1}{2\omega_0}\exp(-|\tau|\omega_0)
        \end{equation}
        while at finite temperature it is periodic with period $\tau_\beta$, and it is given by
        \begin{equation}
            G_E(\tau;\beta) = \frac{\hbar}{m}\,\frac{1}{2\omega_0}
            \frac{\cosh\!\left(\omega_0(\tau_\beta/2 - |\tau|)\right)}{\sinh\!\left(\omega_0\tau_\beta/2\right)},
            \qquad 0 \leq \tau \leq \tau_\beta
        \end{equation}
        \begin{solution}
            For completeness, we shall first show that the propagator is indeed the Green's function of the differential operator given. 
    
            First, note that the Euclidean action can be written in the form
            \begin{equation}
                S_E[q_E]  = \int d\tau \,q_E\frac{m}{2}\left(-\frac{d^2}{d\tau^2}+\omega_0^2\right)q_E.
            \end{equation}
            Using the identity, 
            \begin{equation}
                0= \oint\mathcal{D}q_E\,\frac{\delta}{\delta q_E(\tau_2)}\left(q_E(\tau_1)e^{-S_E[q_E]/\hbar    }\right),
            \end{equation}
            we find
            \begin{equation}
                Z\delta(\tau_1-\tau_2) = \frac{m}{2\hbar}\Tr{e^{-\beta H/\hbar  } q_E(\tau_1)\left(-\frac{d^2}{d\tau^2}+\omega_0^2\right)q_E(\tau_2)},
            \end{equation}
            which implies
            \begin{equation}
                \left(-\frac{d^2}{d\tau^2}+\omega_0^2\right)\Tr{\boldsymbol{\rho}(\beta)\mathbf{Q}(-i\tau_2)\mathbf{Q}(-i\tau_1)} = \frac{\hbar}{m}\delta(\tau_1-\tau_2).
            \end{equation}
            The solution to the Green's function is given by 
            \begin{equation}
                G(\tau;\beta) = \begin{cases}
                    Ae^{+\omega_0\tau} + Be^{-\omega_0\tau}, &\tau<0,\\
                    Ce^{+\omega_0\tau} + De^{-\omega_0\tau}, &\tau>0.
                \end{cases}
            \end{equation}
            For the Green's function to satisfy the periodic boundary condition, we require
            \begin{equation}
                Ce^{+\omega_0\tau} + De^{-\omega_0\tau}  = Ce^{+\omega_0(\tau+\tau_\beta)} + De^{-\omega_0(\tau+\tau_\beta)},
            \end{equation}
            and
            \begin{equation}
                Ae^{+\omega_0\tau} + Be^{-\omega_0\tau}  = Ae^{+\omega_0(\tau-\tau_\beta)} + Be^{-\omega_0(\tau-\tau_\beta)}
            \end{equation}
            for all $\tau$. Therefore, we find
            \begin{equation}
                B = A \frac{1-e^{-\omega_0\tau_\beta}}{e^{\omega_0\tau_\beta}-1}, \quad \text{and} \quad C = D \frac{1-e^{-\omega_0\tau_\beta}}{e^{\omega_0\tau_\beta}-1}.
            \end{equation}
            The matching condition at $\tau=0$ implies
            \begin{equation}
                A + D = \frac{\hbar}{2m\omega_0 }.
            \end{equation}
            For $T\to 0$, $\tau_\beta \to \infty$, we find
            \begin{equation}
                G_E (\tau) = \frac{\hbar}{m}\,\frac{1}{2\omega_0}\exp(-|\tau|\omega_0).
            \end{equation}
            For finite $\tau_\beta$, we find
            \begin{equation}
                G_E(\tau;\beta) = \frac{\hbar}{m}\,\frac{1}{2\omega_0}
            \frac{\cosh\!\left(\omega_0(\tau_\beta/2 - |\tau|)\right)}{\sinh\!\left(\omega_0\tau_\beta/2\right)},
            \qquad 0 \leq \tau \leq \tau_\beta.
            \end{equation}
        \end{solution}
        
    
        \item Compute the Fourier transform $\hat{G}_E(k)$ of $G_E(\tau)$.
        \begin{solution}
            Go back to the differential equation satisfied by $G_E$, we can immediately 
obtain its Fourier transform, when $\tau_\beta \to \infty$:
            \begin{equation}
                \hat{G}_E(k) = \frac{\hbar}{m(k^2+\omega_0^2)}.
            \end{equation}
        \end{solution}
    
    \end{enumerate}
    
    We now go back to the observables at real time $t$.
    
    \begin{enumerate}[label=\alph*), resume]
    
        \item Performing a Wick rotation $\tau = it$, give the explicit form of the Green function $G(t_2, t_1)$. Can you explain why it must be an even function $G(t)$ of $t = t_2 - t_1$?
        \begin{solution}
            Now, we go back to real time. Note that the Euclidean theory with $T=0$ is equivalent to the vacuum expectation value. Therefore, after wick rotation, we find
            \begin{equation}
                G(t_2,t_1) = \frac{\hbar}{m}\frac{1}{2\omega_0}\exp(-i\omega_0|t|)
            \end{equation}
            with $t=t_2-t_1$. It is an even function due to the fact that the ground state of a harmonic oscillator is even under time parity.
        \end{solution}
    
        \item Show that its Fourier transform $\hat{G}(\omega)$ can be written as
        \begin{equation}
            \hat{G}(\omega) = \lim_{\epsilon \to 0^+} \frac{\hbar}{m}\,\frac{i}{\omega^2 - \omega_0^2 + i\epsilon}
        \end{equation}
        \begin{solution}
            We can write the Fourier transform as
    \begin{equation}
        \hat{G}(\omega) = \frac{\hbar}{m}\frac{1}{2\omega_0}
        \int_{-\infty}^{\infty} dt\, e^{i\omega t}e^{-i\omega_0|t|}
    \end{equation}
    We split the integral into $t>0$ and $t<0$ regions, introducing $\epsilon\to 0^+$ in each to ensure convergence:
    \begin{align}
        \hat{G}(\omega) 
        &= \frac{\hbar}{m}\frac{1}{2\omega_0}\left[
           \int_0^\infty dt\, e^{i(\omega-\omega_0+i\epsilon)t} 
           + \int_{-\infty}^0 dt\, e^{i(\omega+\omega_0-i\epsilon)t}
           \right]
    \end{align}
    The $i\epsilon$ in the first integral ensures the integrand vanishes as $t\to+\infty$, and the $-i\epsilon$ in the second ensures it vanishes as $t\to-\infty$. Evaluating both Gaussian integrals exactly:
    \begin{align}
        \hat{G}(\omega) 
        &= \frac{\hbar}{m}\frac{1}{2\omega_0}\left[
           \frac{i}{\omega-\omega_0+i\epsilon} 
           - \frac{i}{\omega+\omega_0-i\epsilon}
           \right] \notag\\
        &= \frac{\hbar}{m}\frac{1}{2\omega_0}\cdot
           \frac{2i\omega_0}{\omega^2-\omega_0^2+i\epsilon}
    \end{align}
    where in the last step we absorbed the $\pm i\epsilon\omega_0$ cross terms into a single $i\epsilon$. This gives exactly
    \begin{equation}
        \hat{G}(\omega) = \lim_{\epsilon\to 0^+}\frac{\hbar}{m}
        \cdot\frac{i}{\omega^2 - \omega_0^2 + i\epsilon}
    \end{equation}
    No approximation has been made — the $(\omega+\omega_0)$ factors cancel exactly rather than requiring any near-pole expansion.
            
        \end{solution}
    
    \end{enumerate}
    
    Going back to the Euclidean Green's function at finite temperature and Euclidean time $G_E(\tau;\beta)$:
    
    \begin{enumerate}[label=\alph*), resume]
    
        \item Express $G_E(\tau;\beta)$ as an infinite sum of zero temperature Green's functions $G_E(\tau + n\tau_\beta)$.
        \begin{solution}
    The finite-temperature Green's function must be periodic with period $\tau_\beta$. 
    A natural way to construct a periodic function from a non-periodic one is the 
    \textit{method of images} — sum the zero-temperature Green's function over all 
    periodic images:
    \begin{equation}
        G_E(\tau;\beta) = \sum_{n=-\infty}^{\infty} G_E(\tau + n\tau_\beta)
        = \frac{\hbar}{m}\frac{1}{2\omega_0}
        \sum_{n=-\infty}^{\infty} e^{-\omega_0|\tau + n\tau_\beta|}
    \end{equation}
    To verify this equals the closed-form expression, restrict to $0\leq\tau\leq\tau_\beta$ 
    so that only the $n\leq 0$ terms contribute positively and $n>0$ negatively:
    \begin{align}
        \sum_{n=-\infty}^{\infty} e^{-\omega_0|\tau+n\tau_\beta|}
        &= \sum_{n=0}^{\infty} e^{-\omega_0(\tau+n\tau_\beta)}
         + \sum_{n=1}^{\infty} e^{\omega_0(\tau-n\tau_\beta)} \notag\\
        &= e^{-\omega_0\tau}\sum_{n=0}^{\infty} e^{-n\omega_0\tau_\beta}
         + e^{\omega_0\tau}\sum_{n=1}^{\infty} e^{-n\omega_0\tau_\beta} \notag\\
        &= \frac{e^{-\omega_0\tau}}{1-e^{-\omega_0\tau_\beta}}
         + \frac{e^{\omega_0\tau}\,e^{-\omega_0\tau_\beta}}{1-e^{-\omega_0\tau_\beta}} \notag\\
        &= \frac{e^{-\omega_0\tau} + e^{\omega_0(\tau-\tau_\beta)}}{1-e^{-\omega_0\tau_\beta}}\notag\\
         &= \frac{2\cosh(\omega_0(\tau_\beta/2-\tau))}{2\sinh(\omega_0\tau_\beta/2)}
    \end{align}
    Therefore:
    \begin{equation}
        G_E(\tau;\beta) = \frac{\hbar}{m}\frac{1}{2\omega_0}
        \frac{\cosh\!\left(\omega_0(\tau_\beta/2-|\tau|)\right)}
             {\sinh\!\left(\omega_0\tau_\beta/2\right)}
    \end{equation}
    which is exactly the closed-form result, confirming the equivalence.
    \end{solution}
    
        \item Using the Poisson summation formula
        \begin{equation}
            \sum_{n=-\infty}^{\infty} f(x + an) = \frac{1}{a}\sum_{k=-\infty}^{\infty} e^{2\pi i k x/a}\,\hat{f}(2\pi k/a)
        \end{equation}
        show that you can rewrite $G_E(\tau;\beta)$ as a sum over frequencies
        \begin{equation}
            G_E(\tau;\beta) = \frac{\hbar}{m}\,\frac{1}{\tau_\beta}
            \sum_{k \in \mathbb{Z}} \frac{e^{i\bar{\omega}_k \tau}}{\omega_k^2 + \omega_0^2}
        \end{equation}
        with $\bar{\omega}_k = -i\omega_k$ given by
        \begin{equation}
            \bar{\omega}_k = \frac{2\pi}{\tau_\beta}\,k, \qquad k \in \mathbb{Z}
        \end{equation}
        \begin{solution}
    Set $x = \tau$, $a = \tau_\beta$ in the Poisson summation formula so that 
    $f(x+an) = e^{-\omega_0|\tau+n\tau_\beta|}$. Using the Fourier transform result 
    $\hat{f}(p) = 2\omega_0/(\omega_0^2+p^2)$ directly:
    \begin{align}
        G_E(\tau;\beta) 
        &= \frac{\hbar}{m}\frac{1}{2\omega_0}
           \sum_{n=-\infty}^{\infty} e^{-\omega_0|\tau+n\tau_\beta|} \notag\\
        &= \frac{\hbar}{m}\frac{1}{2\omega_0}\cdot\frac{1}{\tau_\beta}
           \sum_{k\in\mathbb{Z}} e^{2\pi ik\tau/\tau_\beta}
           \frac{2\omega_0}{\omega_0^2 + (2\pi k/\tau_\beta)^2} \notag\\
        &= \frac{\hbar}{m}\frac{1}{\tau_\beta}
           \sum_{k\in\mathbb{Z}}\frac{e^{i\bar{\omega}_k\tau}}{\omega_0^2+\bar{\omega}_k^2}
    \end{align}
    Identifying $\bar{\omega}_k = 2\pi k/\tau_\beta$ and writing 
    $\bar{\omega}_k^2 = -\omega_k^2$ gives the result.
    \end{solution}
    
        \item The $\omega_k = i\bar{\omega}_k$ can be considered imaginary frequencies. They are called the \textit{Matsubara frequencies} and play an important role in thermal quantum field theory. How do they behave in the limit of zero and infinite temperature?
    \begin{solution}
        As $T \to 0^+$, $\omega_k \to i0^+$. As $T \to \infty$, $\omega_k \to i \infty$.
    \end{solution}
    \end{enumerate}
    
    
    Consider the Green's function at real time
    \begin{equation}
        G(t_1, t_2; \beta) = \mathrm{Tr}\!\left(\rho(\beta)\,T(\mathbf{Q}(t_1)\mathbf{Q}(t_2))\right)
    \end{equation}
    
    \begin{enumerate}[label=\alph*), resume]
    
        \item Express $G(t_1, t_2;\beta)$ explicitly as a function of the time $t = t_2 - t_1$ and of the temperature $T$ or the Boltzmann factor $\beta$.
        \begin{solution}
            Performing a Wick rotation, we find
            \begin{equation}
                G(t;\beta) = \frac{\hbar}{m}\,\frac{i}{2\omega_0}
            \frac{\cos\!\left(\omega_0(\hbar\beta/2 - |t|)\right)}{\sin\!\left(\omega_0\hbar\beta/2\right)},
            \qquad 0 \leq t \leq \hbar\beta.
            \end{equation}
        \end{solution}
    
        \item Express its Fourier transform $\hat{G}(\omega;\beta)$ as a sum over Matsubara frequencies.
        \begin{solution}
    Using the result from part (i):
    \begin{equation}
        G(t;\beta) = \frac{1}{m\beta}\sum_{k\in\mathbb{Z}}
        \frac{e^{i\omega_k t}}{\omega_0^2-\omega_k^2}
    \end{equation}
    Taking the Fourier transform term by term:
    \begin{equation}
        \hat{G}(\omega;\beta) = \int_{-\infty}^{\infty} dt\, e^{i\omega t}\, G(t;\beta)
        = \frac{1}{m\beta}\sum_{k\in\mathbb{Z}}\frac{1}{\omega_0^2-\omega_k^2}
        \int_{-\infty}^{\infty} dt\, e^{i(\omega+\omega_k)t}
    \end{equation}
    Using $\int_{-\infty}^{\infty} dt\, e^{i(\omega+\omega_k)t} = 2\pi\delta(\omega+\omega_k)$
    as a formal identity in the sense of analytic continuation:
    \begin{equation}
        \hat{G}(\omega;\beta) = \frac{2\pi}{m\beta}
        \sum_{k\in\mathbb{Z}}\frac{\delta(\omega+\omega_k)}{\omega_0^2-\omega_k^2}
    \end{equation}
    which is a sum of poles located at $\omega = -\omega_k$ in the complex 
    frequency plane, i.e.\ at the Matsubara frequencies.
\end{solution}
        
    
    \end{enumerate}
    
    